\subsection{分层求解与最优性核验}

附件一电价严格为正，且计划未利用量不产生收益。因此，任意 $q_t>h_t$ 的可行解均可通过令 $q_t=h_t$ 降低费用，第一问求解时可以消去未利用计划变量，将普通购电与实际取电合并。保留充放电互斥变量后，模型包含 $144$ 个二元变量、$721$ 个连续变量、$288$ 个平衡及状态等式和 $288$ 个互斥线性不等式，另以变量界处理初末电量及设备边界。

使用 Python 3.9.21、SciPy 1.13.1 提供的 HiGHS 接口\cite{scipy-milp}求解，设置整数相对最优性间隙为 $10^{-10}$，每个阶段的时间上限为 $120$ 秒。首先最小化购电费，获得最优费用 $C_1^*$；随后在费用不超过 $C_1^*+10^{-7}$ 元的条件下最小化 $\sum_t(c_t+d_t)$；最后，在前两级容差内最小化日循环购电总变差
\begin{equation}
  \sum_{t=0}^{T-1}|q_t-q_{(t-1)\bmod T}|.
  \label{eq:q1-variation}
\end{equation}
后两级目标用于选择经济上近等价的方案，不能替代最低费用目标，也不保证所有时段的购电曲线平滑。最终固定已求得的充放电模式，以 $10^{-9}$ 的线性规划可行性容差重新求解连续电量，消除整数求解容差下的微小负值。

主费用整数模型报告最优解，最优性间隙为 $0$。其最优费用与线性松弛下界满足
\[
  C_1^*=C_{\rm LP}=35126.848589290\ \text{元}.
\]
线性松弛的原始目标与对偶目标一致，且其解满足充放电互斥。最终所选方案比该下界高约 $10^{-7}$ 元，处于预设的择优容差内。因此，下面的费用在分币精度上达到该模型的全局最优。

\subsection{购电与充放电结果}

主方案全天购电量为 $59482.699\ \mathrm{kWh}$，购电费为 $35126.85$ 元。指定时段购电结果见表~\ref{tab:q1-purchase}，分时段储能运行结果见表~\ref{tab:q1-storage}。完整 $144$ 时段策略保存于 \texttt{result1.xlsx}，并另存包含输入、SOC、损耗与费用的逐时段明细。

\input{../output/tables/q1_purchase}
\input{../output/tables/q1_storage}

全天母线侧充电量为 $20740.666\ \mathrm{kWh}$，放电量为 $16799.940\ \mathrm{kWh}$，初末内部储电量均为 $6000\ \mathrm{kWh}$。按式~\eqref{eq:aggregate-balance}，总负荷 $111024.808\ \mathrm{kWh}$、总光伏 $55482.836\ \mathrm{kWh}$ 与转换损耗 $3940.727\ \mathrm{kWh}$ 共同确定实际购电量；本方案光伏弃电量为 $0$。

图~\ref{fig:q1-dispatch} 展示了日内调度。凌晨与 $5{:}40$ 附近的较低价格支持储能充电，随后电池在早间负荷上升阶段放电。中午既吸收富余光伏，也在部分低价时段从外网补充电量，使储能为晚间高价时段保留资源。晚间放电后电量触及下限，$22{:}00$ 后再利用低价购电恢复日末电量。

\begin{figure}[htbp]
  \centering
  \includegraphics[width=\textwidth]{q1_dispatch.pdf}
  \caption{确定性调度的日内对应关系。a：供需与计划购电；b：放电为正、充电为负的储能功率；c：内部储电量及上下限虚线；d：电价与LP净负荷边际成本。功率与电价按10分钟区间绘制。}
  \label{fig:q1-dispatch}
\end{figure}

外网购电曲线在部分低价时段出现较高功率，其峰值为 $8458.827\ \mathrm{kW}$。这不违反题面给出的 $5000\ \mathrm{kW}$ 上限，因为后者限制储能充放电功率，而非外网总取电功率。因此，本模型的经济节省不能直接解释为外网峰值功率下降。图中的净负荷边际成本取自 LP 对偶值，用于解释给定约束状态下增加净负荷的局部代价，不作为任意幅度扰动下的固定价格。

\subsection{基线比较与独立校验}

设置无储能与分位电价规则两种基线。规则基线优先吸收剩余光伏，在已知日价格不高于 $25\%$ 分位数时充电、不低于 $75\%$ 分位数时放电，并加入期末电量可达性修正；该规则未依据优化结果调整阈值。三种策略使用相同输入与日初、日末电量条件，结果见表~\ref{tab:q1-baselines}。

\input{../output/tables/q1_baselines}

相较于无储能方案，优化方案节省 $12925.20$ 元，节费率为 $26.90\%$；相较于规则方案，购电费由 $39573.12$ 元降至 $35126.85$ 元。无储能时有 $6247.963\ \mathrm{kWh}$ 光伏无法利用，优化方案通过跨时段存储实现全部消纳。规则方案仍存在 $3656.782\ \mathrm{kWh}$ 弃光，说明仅按局部价格决定充放电不足以协调未来光伏与容量需求。

独立程序不使用优化矩阵，直接根据输出电量重新核算能量平衡、SOC 转移、功率上限、互斥和费用。主方案的能量平衡及状态转移残差均小于 $10^{-8}\ \mathrm{kWh}$，储电量保持在 $1200$--$10800\ \mathrm{kWh}$ 内，初末电量误差为 $0$。CSV 回读与官方工作簿回读检查同时通过，工作表名称、非结果单元格内容及格式保持不变。

\subsection{参数与时间解释的影响}

对容量、充放电功率及效率进行单因素重新优化，结果如图~\ref{fig:q1-sensitivity}。容量改变时同步保持 SOC 下限为额定容量的 $10\%$、上限为 $90\%$、初末电量为 $50\%$；因此，该组结果代表整套储能规格变化，不能解释为固定初始电量下的纯容量边际收益。

\begin{figure}[htbp]
  \centering
  \includegraphics[width=\textwidth]{q1_sensitivity.pdf}
  \caption{设备参数对最优费用的影响。a--c分别改变容量、功率上限与单程效率；各点均重新优化，实心点及竖虚线对应主设定。纵轴为局部费用范围，连线不表示统计置信区间。}
  \label{fig:q1-sensitivity}
\end{figure}

额定容量由 $6000$ 增至 $18000\ \mathrm{kWh}$ 时，最优费用由 $40020.40$ 降至 $32093.06$ 元；功率上限由 $2500$ 增至 $5000\ \mathrm{kW}$ 时，费用由 $36506.41$ 降至 $35126.85$ 元，再增至 $7500\ \mathrm{kW}$ 时为 $35027.13$ 元。这表明，在本日输入下，提高功率上限的边际节费趋于减弱，但该观察不能替代包含设备投资的容量选择分析。

若将题面的 $90\%$ 解释为往返效率，并取 $\eta_c=\eta_d=\sqrt{0.9}$，费用为 $33801.48$ 元，说明效率解释需要在后续各问中保持一致。光伏输入分别缩放为原值的 $0.8$ 和 $1.2$ 倍并重新优化时，费用分别为 $42268.99$ 和 $30873.96$ 元；这些属于确定性输入敏感性，不能视为既定计划应对真实预测误差的表现。

将附件时间改作时段终点，并在同样初末电量条件下重新优化，费用为 $35126.95$ 元，与主解释相差约 $0.10$ 元。虽然全天费用接近，指定时段的映射并不相同，因此仍保存两套完整策略及逐项对应关系。主结果以模板时段起点解释为准，为后续各问提供统一的时间、设备及费用核算接口。
